% GENERATED FILE. Edit canonical lessons, then run npm run generate:books.
% canonical-source-hash: fnv1a64:fc74953cdc29470d
% canonical-lessons: FR-C276-epoque, FR-C276-decennie, FR-C276-millenaire, FR-C276-ere, FR-C276-duree

\chapter{Les époques}
\label{ch:fr-les-epoques}
\hlchaptermodality{276}

\begin{chapteropening}
\textbf{By the end of this chapter:} \emph{I can talk about stretches of time: an era, a decade, a millennium, an age and a duration.}

Complete the last lesson of chapter 276: I can talk about stretches of time: an era, a decade, a millennium, an age and a duration.
\end{chapteropening}

\section[l'époque]{l'époque — the time, the era, the period}

\label{lesson:FR-C276-epoque}



\begin{quote}
Before the new one: say the French for the hitch, then the French for the deadline.
\end{quote}

\subsection*{You'll want to know: l'époque}
\textbf{l'époque} — \textquotedblleft{}the time, the era, the period\textquotedblright{}.

\begin{quote}
À l'époque, on n'avait pas de téléphone portable.
\end{quote}

\begin{quote}
\emph{Back then, people did not have mobile phones.}
\end{quote}

\emph{À l'époque} — back then; \emph{à cette époque-là} — at that time.

\begin{cousinweb}
\emph{époque} is a learned borrowing from Greek \emph{epokh\'{\={e}}} \textquotedblleft{}a halt, a fixed point in time\textquotedblright{}, which reached French through Medieval Latin (EtymDB lists the French word without a parent).
\end{cousinweb}

\begin{grammarlens}[title={What you've built}]
One more word for telling, explaining and describing at B1.
\end{grammarlens}

\subsection*{Your turn}
\begin{itemize}
\raggedright
  \item \emph{Say it:} \emph{l'époque}
  \item \emph{Say it:} \emph{l'époque}, once more
  \item \emph{Say it:} a sentence of your own with \emph{l'époque}
  \item \emph{Recall:} say the French for the hitch, then the French for the deadline, then say \emph{l'époque} again
\end{itemize}

\begin{tcolorbox}[breakable,colback=teal!4,colframe=teal!35!black,title={Before you move on}]
What does l'époque mean? (\textquotedblleft{}The time, the era, the period\textquotedblright{}.) Say it once more.
\end{tcolorbox}
\section[la décennie]{la décennie — the decade}

\label{lesson:FR-C276-decennie}



\begin{quote}
Before the new one: say the French for the deadline, then the French for the time.
\end{quote}

\subsection*{You'll want to know: la décennie}
\textbf{la décennie} — \textquotedblleft{}the decade\textquotedblright{}.

\begin{quote}
La ville a beaucoup changé en une décennie.
\end{quote}

\begin{quote}
\emph{The town changed a lot in a decade.}
\end{quote}

Ten years. \emph{La décade} is also heard in this sense, though strictly it means ten days.

\begin{cousinweb}
\emph{décennie} is a learned borrowing from Latin \emph{decennium} \textquotedblleft{}a period of ten years\textquotedblright{} (EtymDB), from \emph{decem} \textquotedblleft{}ten\textquotedblright{} and \emph{annus} \textquotedblleft{}year\textquotedblright{}.
\end{cousinweb}

\begin{grammarlens}[title={What you've built}]
One more word for telling, explaining and describing at B1.
\end{grammarlens}

\subsection*{Your turn}
\begin{itemize}
\raggedright
  \item \emph{Say it:} \emph{la décennie}
  \item \emph{Say it:} \emph{la décennie}, once more
  \item \emph{Say it:} a sentence of your own with \emph{la décennie}
  \item \emph{Recall:} say the French for the deadline, then the French for the time, then say \emph{la décennie} again
\end{itemize}

\begin{tcolorbox}[breakable,colback=teal!4,colframe=teal!35!black,title={Before you move on}]
What does la décennie mean? (\textquotedblleft{}The decade\textquotedblright{}.) Say it once more.
\end{tcolorbox}
\section[le millénaire]{le millénaire — the millennium}

\label{lesson:FR-C276-millenaire}



\begin{quote}
Before the new one: say the French for the time, then the French for the decade.
\end{quote}

\subsection*{You'll want to know: le millénaire}
\textbf{le millénaire} — \textquotedblleft{}the millennium\textquotedblright{}.

\begin{quote}
Cette église a presque un millénaire.
\end{quote}

\begin{quote}
\emph{This church is almost a thousand years old.}
\end{quote}

\emph{Le troisième millénaire} — the third millennium.

\begin{cousinweb}
\emph{millénaire} is a learned borrowing from Latin \emph{millēnārius} \textquotedblleft{}containing a thousand\textquotedblright{} (EtymDB), built on \emph{mīlle} \textquotedblleft{}thousand\textquotedblright{}.
\end{cousinweb}

\begin{grammarlens}[title={What you've built}]
One more word for telling, explaining and describing at B1.
\end{grammarlens}

\subsection*{Your turn}
\begin{itemize}
\raggedright
  \item \emph{Say it:} \emph{le millénaire}
  \item \emph{Say it:} \emph{le millénaire}, once more
  \item \emph{Say it:} a sentence of your own with \emph{le millénaire}
  \item \emph{Recall:} say the French for the time, then the French for the decade, then say \emph{le millénaire} again
\end{itemize}

\begin{tcolorbox}[breakable,colback=teal!4,colframe=teal!35!black,title={Before you move on}]
What does le millénaire mean? (\textquotedblleft{}The millennium\textquotedblright{}.) Say it once more.
\end{tcolorbox}
\section[l'ère]{l'ère — the era, the age}

\label{lesson:FR-C276-ere}



\begin{quote}
Before the new one: say the French for the decade, then the French for the millennium.
\end{quote}

\subsection*{You'll want to know: l'ère}
\textbf{l'ère} — \textquotedblleft{}the era, the age\textquotedblright{}.

\begin{quote}
Nous vivons dans l'ère d'internet.
\end{quote}

\begin{quote}
\emph{We live in the internet age.}
\end{quote}

Feminine: \emph{une nouvelle ère}, a new era.

\begin{cousinweb}
\emph{ère} was borrowed from Late Latin \emph{aera} \textquotedblleft{}a counted number, an era\textquotedblright{} (EtymDB).
\end{cousinweb}

\begin{grammarlens}[title={What you've built}]
One more word for telling, explaining and describing at B1.
\end{grammarlens}

\subsection*{Your turn}
\begin{itemize}
\raggedright
  \item \emph{Say it:} \emph{l'ère}
  \item \emph{Say it:} \emph{l'ère}, once more
  \item \emph{Say it:} a sentence of your own with \emph{l'ère}
  \item \emph{Recall:} say the French for the decade, then the French for the millennium, then say \emph{l'ère} again
\end{itemize}

\begin{tcolorbox}[breakable,colback=teal!4,colframe=teal!35!black,title={Before you move on}]
What does l'ère mean? (\textquotedblleft{}The era, the age\textquotedblright{}.) Say it once more.
\end{tcolorbox}
\section[la durée]{la durée — the length, the duration}

\label{lesson:FR-C276-duree}



\begin{quote}
Before the new one: say the French for the millennium, then the French for the era.
\end{quote}

\subsection*{You'll want to know: la durée}
\textbf{la durée} — \textquotedblleft{}the length, the duration\textquotedblright{}.

\begin{quote}
La durée du film est de deux heures.
\end{quote}

\begin{quote}
\emph{The film is two hours long.}
\end{quote}

\emph{Pour une durée de six mois} — for a period of six months.

\begin{cousinweb}
\emph{durée} is built on \emph{durer} \textquotedblleft{}to last\textquotedblright{}, which is inherited from Latin \emph{dūrāre} \textquotedblleft{}to hold out, to last\textquotedblright{} (EtymDB).
\end{cousinweb}

\begin{grammarlens}[title={What you've built}]
One more word for telling, explaining and describing at B1.
\end{grammarlens}

\subsection*{Your turn}
\begin{itemize}
\raggedright
  \item \emph{Say it:} \emph{la durée}
  \item \emph{Say it:} \emph{la durée}, once more
  \item \emph{Say it:} a sentence of your own with \emph{la durée}
  \item \emph{Recall:} say the French for the millennium, then the French for the era, then say \emph{la durée} again
\end{itemize}

\begin{tcolorbox}[breakable,colback=teal!4,colframe=teal!35!black,title={Before you move on}]
What does la durée mean? (\textquotedblleft{}The length, the duration\textquotedblright{}.) Say it once more.
\end{tcolorbox}
